% Chapter 11: Conclusions

\chapter{Conclusions}

This thesis presented \textbf{SatelliteCV-Paraguay}, a comprehensive open-source
Python toolkit that unifies six peer-reviewed papers across remote sensing,
agriculture, climate, conservation, and social science. The toolkit
demonstrates that foundation models pre-trained on global Earth observation
data can be successfully transferred to Paraguay-specific applications.

\section{Contributions}

The main contributions of this thesis are:

\begin{enumerate}
    \item \textbf{Open-source Python package} (\texttt{satellite-paraguay})
    with 6 paper pipelines, baseline implementations, FastAPI endpoint,
    Streamlit dashboard, and comprehensive documentation.

    \item \textbf{Six peer-reviewed papers} covering:
    \begin{itemize}
        \item P0011 Yvutu — Chaco deforestation (Remote Sensing of Environment)
        \item P0100 Yvyra — Carbon credit verification (Nature Climate Change)
        \item P0025 Yrupe — Soybean yield (Computers and Electronics in Agriculture)
        \item P0012 Yvy — Indigenous territory (World Development, CARE-compliant)
        \item P0026 Kai — Wildlife poaching (Conservation Biology)
        \item P0035 Tatakua — Air quality (Atmospheric Environment)
    \end{itemize}

    \item \textbf{Comprehensive data catalog} of 14+ datasets, all open-source.

    \item \textbf{Best practices} for indigenous data sovereignty (CARE Principles).

    \item \textbf{Reproducibility framework} with random seeds, environment
    capture, MLflow tracking, and DVC versioning.

    \item \textbf{Cross-paper insights} that emerge from the unified platform.
\end{enumerate}

\section{Future Work}

\begin{itemize}
    \item Multi-modal extensions (GPT-4V for place name recognition)
    \item Self-supervised pretraining on Paraguay-specific tiles
    \item Federated learning across Paraguayan institutions
    \item Edge AI deployment (Jetson for forest monitoring)
    \item Mobile app for community-driven mapping
    \item Real-time dashboard with live Sentinel-2 ingest
    \item Extension to Bolivia + Argentina (Gran Chaco regional analysis)
\end{itemize}

\section{Final Remarks}

Paraguay has unique natural and cultural resources that deserve world-class
scientific attention. This thesis provides a foundation for continued
Earth observation research in Paraguay and serves as a template for similar
efforts in other Latin American countries.

The complete codebase, datasets, and documentation are released as open
source to support future researchers. We invite collaboration from the
international scientific community.

\vspace{2cm}
\begin{center}
\textit{Para el Paraguay que viene.}
\end{center}
